\subsection{Silver labels and annotation status}

The calibration and evaluation of the knowledge-extraction pipeline does not rely on gold labels annotated by a domain expert. The reference labels, against which the pipeline is calibrated and evaluated, are silver labels \pinktodo{Add Zheng et al. 2023 here, after verifying that the paper is good enough.}: the per-case finding sets are generated by prompting a strong language model (Claude Opus 4.7, developed by Anthropic), \pinktodo{decide if we want to regenerate labels again at temperature zero, to do the same thing Zheng et al. do in their paper) temperature zero, to maximize reproducibility.} The resulting labels therefore represent a strong-LLM-judge approximation of which clinical findings should be extracted from each text chunk. 

\pinktodo{Add temperature inside the silver label cache, and regenerate silver labels by Opus for the ILP-selected 15 papers again with temperature zero.}

\pinktodo{The silver finding cases from Opus are not run on the 15 correct papers anyway, we have to regenerate the silver-findings.}

Silver labels are useful because they scale: \pinktodo{Change the numbers here.} across 273 calibration cases, 1\,596 silver findings and a 300-claim pair evaluation dataset, consisting of claims supporting, contradicting or unrelated to each other were generated; producing the same number of expert-curated annotations would not be within the scope of this thesis. The silver-labels make it possible to calibrate the thresholds and evaluate the pipeline. 

In place of expert annotation, the thesis uses two automatically generated, silver-quality artifacts (Zheng et al., 2023), each with a distinct role: a silver finding set for threshold and hyperparameter calibration, and a separate dataset consisting of 300 claim pairs for the evaluation of the pipeline. 

The silver-finding set is produced by Claude Opus 4.7, prompted with a silver-generation prompt at temperature zero. The zero temperature allows for the set to be reproducible. Across 273 cases, the result is a set of 1\,596 findings. These findings are used for threshold and hyperparameter calibration of the knowledge-extraction pipeline.

The pipeline is then evaluated on a 300-claim pair dataset, spanning a range of histopathology topics, each labelled supporting, contradicting, or unrelated. At this stage, the thresholds are not calibrated again, and the sole purpose of this silver-labelled dataset is the evaluation. It targets \pinktodo{all stages/relation stage?} of the pipeline, which, in the end, generate claims from text chunks that come from the document-extraction pipeline, and decide whether the claims support or contradict each other, or are unrelated. 

Neither of the datasets are treated as definitive ground truth. Since the calibration depends on the silver-finding dataset, the thresholds inherit the biases of Claude Opus 4.7. Since nothing is tuned on a different silver-dataset consisting of 300 claim-pairs, the threat of overfitting is eliminated when calculating F1-scores for pipeline evaluation. However, similar to the threshold calibration case on the silver-finding set, evaluation also inherits the biases of the same generator model. These claim pairs might also have cleaner relations among each other: supporting claims might be very similar, contradicting claims very sharply opposing, and unrelated claims more plainly unrelated. Therefore, generalization from 300 claim pairs to a corpus-scale output is not possible, which will further be discussed in \pinktodo{Chapter 5, Discussion}. 



